\section{Introduction}
\label{sec:intro}

Transformers that process many in-context examples narrow their output distributions: with enough examples of a pattern, the model becomes more ``certain'' about what comes next. A natural question follows: {\bf does this distributional narrowing correspond to geometric simplification in the residual stream?} If so, residual stream geometry could serve as a mechanistic signature of in-context learning progress, with applications to interpretability, uncertainty quantification, and understanding the limits of context utilization.

Recent work in mechanistic interpretability has established that transformers linearly encode belief state geometry in their residual streams~\cite{shai2024transformers}, and that game-playing models develop emergent world representations accessible via linear probes~\cite{nanda2023emergent, li2023emergent}. However, these studies analyze geometry at fixed context positions without tracking how it evolves as context grows. The many-shot regime---where models process tens or hundreds of in-context examples---remains geometrically uncharacterized.

We address this gap by training small decoder-only transformers on sequences from games of varying complexity: a 3-state hidden Markov model (\messthree), Rock-Paper-Scissors with counter, random, and win-stay-lose-shift strategies, and \ttt. For each model, we extract residual stream activations at every context position and measure three geometric properties: effective dimensionality via participation ratio, PCA variance concentration, and linear probe accuracy for latent game states.

Our main finding challenges the many-shot simplification hypothesis. Geometric simplification is not a universal consequence of processing more context. Instead, it tracks the \emph{effective entropy} of the data-generating process at each position:
\begin{itemize}[leftmargin=*,itemsep=0pt,topsep=0pt]
    \item \ttt shows dramatic dimensionality collapse ($\PR$: $7.85 \to 1.26$, Cohen's $d = 19.98$) because the game's state space shrinks as the board fills---a structural property of the task, not a learning phenomenon.
    \item Stationary processes (\messthree, RPS) show approximately constant geometry across context ($\PR \approx 2.0$--$2.3$), stabilizing within the first ${\sim}5$ positions regardless of how much additional context is provided.
    \item Belief state regression quality unexpectedly \emph{decreases} with context position in \messthree ($R^2$: $1.0 \to 0.95$, $p < 0.001$), revealing representational drift rather than progressive refinement.
\end{itemize}

\para{Contributions.} We make the following contributions:
\begin{itemize}[leftmargin=*,itemsep=0pt,topsep=0pt]
    \item We provide the first systematic measurement of residual stream geometry as a function of context position across multiple game-playing tasks, testing the many-shot geometric simplification hypothesis.
    \item We demonstrate that effective dimensionality tracks task-level state-space entropy rather than context accumulation, with $\PR \approx 2$ emerging as a consistent floor for 3-token vocabularies---consistent with the belief state simplex being 2-dimensional.
    \item We identify an unexpected belief state regression degradation with context position, suggesting that small transformers experience representational drift over long sequences.
\end{itemize}

The remainder of this paper is organized as follows. \Secref{sec:related} reviews related work on residual stream geometry and in-context learning. \Secref{sec:method} describes our experimental setup, including data generation, model architecture, and geometric metrics. \Secref{sec:results} presents our findings across all six datasets. \Secref{sec:discussion} interprets the results and discusses limitations. \Secref{sec:conclusion} summarizes our contributions and outlines future directions.
